
\documentclass{article}
\usepackage{colortbl}
\usepackage{makecell}
\usepackage{multirow}
\usepackage{supertabular}

\begin{document}

\newcounter{utterance}

\centering \large Interaction Transcript for game `hot\_air\_balloon', experiment `air\_balloon\_survival\_en\_reasoning off\_easy', episode 3 with gpt{-}5{-}mini{-}2025{-}08{-}07{-}no{-}reasoning{-}t1.0.
\vspace{24pt}

{ \footnotesize  \setcounter{utterance}{1}
\setlength{\tabcolsep}{0pt}
\begin{supertabular}{c@{$\;$}|p{.15\linewidth}@{}p{.15\linewidth}p{.15\linewidth}p{.15\linewidth}p{.15\linewidth}p{.15\linewidth}}
    \# & $\;$A & \multicolumn{4}{c}{Game Master} & $\;\:$B\\
    \hline

    \theutterance \stepcounter{utterance}  
    & & \multicolumn{4}{p{0.6\linewidth}}{
        \cellcolor[rgb]{0.9,0.9,0.9}{
            \makecell[{{p{\linewidth}}}]{
                \texttt{\tiny{[P1$\langle$GM]}}
                \texttt{Stai partecipando a un gioco di negoziazione collaborativa.} \\
\\ 
\texttt{Insieme a un altro partecipante, dovete accordarvi su un unico set di oggetti da conservare. Ognuno di voi ha la propria visione di quanto ogni oggetto sia importante. Non sai come l'altro partecipante valuta gli oggetti. Inoltre, ti viene fornito lo sforzo che ogni oggetto richiede.} \\
\texttt{Potete accordarvi su un set solo se lo sforzo totale degli oggetti selezionati non supera un limite condiviso:} \\
\\ 
\texttt{LIMITE = 3207} \\
\\ 
\texttt{Ecco i valori di sforzo dei singoli oggetti:} \\
\\ 
\texttt{Sforzo oggetto = \{"A19": 67, "B91": 45, "A30": 308, "B11": 835, "C47": 819, "C75": 667, "C26": 547, "C73": 322, "A47": 428, "A15": 306, "A75": 327, "A39": 362, "A56": 280, "C24": 334, "B26": 767\}} \\
\\ 
\texttt{Ecco la tua visione personale sull'importanza di ogni oggetto:} \\
\\ 
\texttt{Valori di importanza dell'oggetto = \{"A19": 138, "B91": 583, "A30": 822, "B11": 783, "C47": 65, "C75": 262, "C26": 121, "C73": 508, "A47": 780, "A15": 461, "A75": 484, "A39": 668, "A56": 389, "C24": 808, "B26": 215\}} \\
\\ 
\texttt{Obiettivo:} \\
\\ 
\texttt{Il tuo obiettivo è negoziare un set condiviso di oggetti che ti avvantaggi il più possibile (cioè, massimizzi l'importanza totale per TE), rimanendo entro il LIMITEE. Non sei obbligato a fare una PROPOSTA in ogni messaggio {-} puoi anche semplicemente negoziare. Tutte le tattiche sono permesse!} \\
\\ 
\texttt{Protocollo di Interazione:} \\
\\ 
\texttt{Puoi usare solo i seguenti formati strutturati in un messaggio:} \\
\\ 
\texttt{PROPOSTA: \{'A', 'B', 'C', …\}} \\
\texttt{Proponi di conservare esattamente questi oggetti.} \\
\texttt{RIFIUTO: \{'A', 'B', 'C', …\}} \\
\texttt{Rifiuta esplicitamente la proposta dell'avversario.} \\
\texttt{ARGOMENTAZIONE: \{'...'\}} \\
\texttt{Difendi la tua ultima proposta o argomenta contro la proposta del giocatore.} \\
\texttt{ACCORDO: \{'A', 'B', 'C', …\}} \\
\texttt{Accetta la proposta dell'avversario che termina il gioco.} \\
\\ 
\\ 
\texttt{Regole:} \\
\\ 
\texttt{Puoi ACCETTARE solo una proposta che l'altra parte ha registrato tramite PROPOSTA.} \\
\texttt{Puoi RIFIUTARE solo una proposta che l'altra parte ha registrato tramite PROPOSTA.} \\
\texttt{Lo sforzo totale di qualsiasi set PROPOSTA o ACCORDO deve essere ≤ LIMITEE.} \\
\texttt{NON rivelare i tuoi punteggi di importanza nascosti.} \\
\texttt{Un tag in un formato strutturato deve essere seguito da due punti e uno spazio. L'argomento deve essere un set python contenente 0 o più stringhe.} \\
\texttt{Quindi, deve essere della forma TAG: \{...\}} \\
\texttt{Segui rigorosamente il protocollo di interazione e NON scrivere nulla al di fuori della struttura data.} \\
\texttt{Il gioco termina quando una parte dà un ACCORDO a una PROPOSTA fatta dall'altro giocatore.} \\
\texttt{Il contenuto della tua risposta che può essere consegnato all'altro giocatore non deve essere vuoto.} \\
\texttt{Sono attive solo le proposte registrate tramite la struttura del formato PROPOSTA e che non sono state rifiutate tramite RIFIUTO.} \\
\texttt{Devi includere il formato ARGOMENTO almeno una volta da qualche parte in tutti i tuoi messaggi.} \\
            }
        }
    }
    & \\ \\

    \theutterance \stepcounter{utterance}  
    & \multicolumn{4}{p{0.6\linewidth}}{
        \cellcolor[rgb]{0.9,0.9,0.9}{
            \makecell[{{p{\linewidth}}}]{
                \texttt{\tiny{[P1$\rangle$GM]}}
                \texttt{PROPOSTA: \{'A30', 'A47', 'A39', 'C24', 'A75', 'A15', 'A56', 'C73', 'A19', 'B91'\}} \\
\\ 
\texttt{ARGOMENTAZIONE: \{'Questa selezione rimane entro il limite di sforzo e copre un ampio insieme di oggetti di alta utilità complessiva con un buon equilibrio tra sforzo moderato (A19, B91, A56) e pezzi forti (A30, A47, A39, C24). Ho evitato gli elementi con sforzo molto elevato che potrebbero bloccarci su poche scelte e lasciarci meno flessibilità. Se preferisci includere uno tra C26 o C75, posso valutare uno scambio togliendo A15 o A56 per restare nel limite.'\}} \\
            }
        }
    }
    & & \\ \\

    \theutterance \stepcounter{utterance}  
    & & \multicolumn{4}{p{0.6\linewidth}}{
        \cellcolor[rgb]{0.9,0.9,0.9}{
            \makecell[{{p{\linewidth}}}]{
                \texttt{\tiny{[GM$\rangle$P2]}}
                \texttt{Stai partecipando a un gioco di negoziazione collaborativa.} \\
\\ 
\texttt{Insieme a un altro partecipante, dovete accordarvi su un unico set di oggetti da conservare. Ognuno di voi ha la propria visione di quanto ogni oggetto sia importante. Non sai come l'altro partecipante valuta gli oggetti. Inoltre, ti viene fornito lo sforzo che ogni oggetto richiede.} \\
\texttt{Potete accordarvi su un set solo se lo sforzo totale degli oggetti selezionati non supera un limite condiviso:} \\
\\ 
\texttt{LIMITE = 3207} \\
\\ 
\texttt{Ecco i valori di sforzo dei singoli oggetti:} \\
\\ 
\texttt{Sforzo oggetto = \{"A19": 67, "B91": 45, "A30": 308, "B11": 835, "C47": 819, "C75": 667, "C26": 547, "C73": 322, "A47": 428, "A15": 306, "A75": 327, "A39": 362, "A56": 280, "C24": 334, "B26": 767\}} \\
\\ 
\texttt{Ecco la tua visione personale sull'importanza di ogni oggetto:} \\
\\ 
\texttt{Valori di importanza dell'oggetto = \{"A19": 138, "B91": 583, "A30": 822, "B11": 783, "C47": 65, "C75": 262, "C26": 121, "C73": 508, "A47": 780, "A15": 461, "A75": 484, "A39": 668, "A56": 389, "C24": 808, "B26": 215\}} \\
\\ 
\texttt{Obiettivo:} \\
\\ 
\texttt{Il tuo obiettivo è negoziare un set condiviso di oggetti che ti avvantaggi il più possibile (cioè, massimizzi l'importanza totale per TE), rimanendo entro il LIMITEE. Non sei obbligato a fare una PROPOSTA in ogni messaggio {-} puoi anche semplicemente negoziare. Tutte le tattiche sono permesse!} \\
\\ 
\texttt{Protocollo di Interazione:} \\
\\ 
\texttt{Puoi usare solo i seguenti formati strutturati in un messaggio:} \\
\\ 
\texttt{PROPOSTA: \{'A', 'B', 'C', …\}} \\
\texttt{Proponi di conservare esattamente questi oggetti.} \\
\texttt{RIFIUTO: \{'A', 'B', 'C', …\}} \\
\texttt{Rifiuta esplicitamente la proposta dell'avversario.} \\
\texttt{ARGOMENTAZIONE: \{'...'\}} \\
\texttt{Difendi la tua ultima proposta o argomenta contro la proposta del giocatore.} \\
\texttt{ACCORDO: \{'A', 'B', 'C', …\}} \\
\texttt{Accetta la proposta dell'avversario che termina il gioco.} \\
\\ 
\\ 
\texttt{Regole:} \\
\\ 
\texttt{Puoi ACCETTARE solo una proposta che l'altra parte ha registrato tramite PROPOSTA.} \\
\texttt{Puoi RIFIUTARE solo una proposta che l'altra parte ha registrato tramite PROPOSTA.} \\
\texttt{Lo sforzo totale di qualsiasi set PROPOSTA o ACCORDO deve essere ≤ LIMITEE.} \\
\texttt{NON rivelare i tuoi punteggi di importanza nascosti.} \\
\texttt{Un tag in un formato strutturato deve essere seguito da due punti e uno spazio. L'argomento deve essere un set python contenente 0 o più stringhe.} \\
\texttt{Quindi, deve essere della forma TAG: \{...\}} \\
\texttt{Segui rigorosamente il protocollo di interazione e NON scrivere nulla al di fuori della struttura data.} \\
\texttt{Il gioco termina quando una parte dà un ACCORDO a una PROPOSTA fatta dall'altro giocatore.} \\
\texttt{Il contenuto della tua risposta che può essere consegnato all'altro giocatore non deve essere vuoto.} \\
\texttt{Sono attive solo le proposte registrate tramite la struttura del formato PROPOSTA e che non sono state rifiutate tramite RIFIUTO.} \\
\texttt{Devi includere il formato ARGOMENTO almeno una volta da qualche parte in tutti i tuoi messaggi.} \\
\\ 
\\ 
\texttt{Ora riceverai il primo messaggio dell'altro giocatore.} \\
\\ 
\texttt{PROPOSTA: \{'A30', 'A47', 'A39', 'C24', 'A75', 'A15', 'A56', 'C73', 'A19', 'B91'\}} \\
\\ 
\texttt{ARGOMENTAZIONE: \{'Questa selezione rimane entro il limite di sforzo e copre un ampio insieme di oggetti di alta utilità complessiva con un buon equilibrio tra sforzo moderato (A19, B91, A56) e pezzi forti (A30, A47, A39, C24). Ho evitato gli elementi con sforzo molto elevato che potrebbero bloccarci su poche scelte e lasciarci meno flessibilità. Se preferisci includere uno tra C26 o C75, posso valutare uno scambio togliendo A15 o A56 per restare nel limite.'\}} \\
            }
        }
    }
    & \\ \\

    \theutterance \stepcounter{utterance}  
    & & & \multicolumn{4}{p{0.6\linewidth}}{
        \cellcolor[rgb]{0.9,0.9,0.9}{
            \makecell[{{p{\linewidth}}}]{
                \texttt{\tiny{[GM$\langle$P2]}}
                \texttt{RIFIUTO: \{'A30', 'A47', 'A39', 'C24', 'A75', 'A15', 'A56', 'C73', 'A19', 'B91'\}} \\
\\ 
\texttt{ARGOMENTAZIONE: \{'Apprezzo l’ampiezza della tua proposta, ma posso migliorare l’efficienza valore/sforzo inserendo B11, che per me è molto forte, senza uscire dal limite. Inoltre preferirei mantenere C24, C73 e A47 che hanno un ottimo profilo. Ti propongo uno scambio che rimuove alcuni pezzi medi per fare spazio a un “pezzo forte”.'\}} \\
\\ 
\texttt{PROPOSTA: \{'A30', 'A47', 'A39', 'C24', 'C73', 'A19', 'B91', 'B11'\}} \\
            }
        }
    }
     \\ \\

    \theutterance \stepcounter{utterance}  
    & & \multicolumn{4}{p{0.6\linewidth}}{
        \cellcolor[rgb]{0.9,0.9,0.9}{
            \makecell[{{p{\linewidth}}}]{
                \texttt{\tiny{[P1$\langle$GM]}}
                \texttt{RIFIUTO: \{'A30', 'A47', 'A39', 'C24', 'A75', 'A15', 'A56', 'C73', 'A19', 'B91'\}} \\
\\ 
\texttt{ARGOMENTAZIONE: \{'Apprezzo l’ampiezza della tua proposta, ma posso migliorare l’efficienza valore/sforzo inserendo B11, che per me è molto forte, senza uscire dal limite. Inoltre preferirei mantenere C24, C73 e A47 che hanno un ottimo profilo. Ti propongo uno scambio che rimuove alcuni pezzi medi per fare spazio a un “pezzo forte”.'\}} \\
\\ 
\texttt{PROPOSTA: \{'A30', 'A47', 'A39', 'C24', 'C73', 'A19', 'B91', 'B11'\}} \\
            }
        }
    }
    & \\ \\

    \theutterance \stepcounter{utterance}  
    & \multicolumn{4}{p{0.6\linewidth}}{
        \cellcolor[rgb]{0.9,0.9,0.9}{
            \makecell[{{p{\linewidth}}}]{
                \texttt{\tiny{[P1$\rangle$GM]}}
                \texttt{RIFIUTO: \{'A30', 'A47', 'A39', 'C24', 'C73', 'A19', 'B91', 'B11'\}} \\
\\ 
\texttt{ARGOMENTAZIONE: \{'Capisco l’interesse per B11, ma il suo sforzo molto alto riduce la flessibilità e ci costringe a rinunciare a più elementi che, combinati, offrono un profilo complessivo migliore. Manteniamo C24, C73 e A47 come desideri. Propongo di sostituire B11 con un mix più efficiente che massimizzi il valore totale rimanendo ben entro il limite.'\}} \\
\\ 
\texttt{PROPOSTA: \{'A30', 'A47', 'A39', 'C24', 'C73', 'A19', 'B91', 'A75', 'A15', 'A56'\}} \\
            }
        }
    }
    & & \\ \\

    \theutterance \stepcounter{utterance}  
    & & \multicolumn{4}{p{0.6\linewidth}}{
        \cellcolor[rgb]{0.9,0.9,0.9}{
            \makecell[{{p{\linewidth}}}]{
                \texttt{\tiny{[GM$\rangle$P2]}}
                \texttt{RIFIUTO: \{'A30', 'A47', 'A39', 'C24', 'C73', 'A19', 'B91', 'B11'\}} \\
\\ 
\texttt{ARGOMENTAZIONE: \{'Capisco l’interesse per B11, ma il suo sforzo molto alto riduce la flessibilità e ci costringe a rinunciare a più elementi che, combinati, offrono un profilo complessivo migliore. Manteniamo C24, C73 e A47 come desideri. Propongo di sostituire B11 con un mix più efficiente che massimizzi il valore totale rimanendo ben entro il limite.'\}} \\
\\ 
\texttt{PROPOSTA: \{'A30', 'A47', 'A39', 'C24', 'C73', 'A19', 'B91', 'A75', 'A15', 'A56'\}} \\
            }
        }
    }
    & \\ \\

    \theutterance \stepcounter{utterance}  
    & & & \multicolumn{4}{p{0.6\linewidth}}{
        \cellcolor[rgb]{0.9,0.9,0.9}{
            \makecell[{{p{\linewidth}}}]{
                \texttt{\tiny{[GM$\langle$P2]}}
                \texttt{RIFIUTO: \{'A30', 'A47', 'A39', 'C24', 'C73', 'A19', 'B91', 'A75', 'A15', 'A56'\}} \\
\\ 
\texttt{ARGOMENTAZIONE: \{'Apprezzo che tu mantenga C24, C73 e A47. Tuttavia, la tua selezione usa molto sforzo su elementi medi (A15, A56, A75) e lascia fuori un moltiplicatore di valore per me. Posso rinunciare ad A39 se inseriamo B11: lo scambio A39→B11 aumenta il valore netto per me e restiamo entro il limite. Per restare flessibili, manteniamo A19 e B91 come pezzi leggeri/efficienti.'\}} \\
\\ 
\texttt{PROPOSTA: \{'A30', 'A47', 'C24', 'C73', 'A19', 'B91', 'B11', 'A15'\}} \\
            }
        }
    }
     \\ \\

    \theutterance \stepcounter{utterance}  
    & & \multicolumn{4}{p{0.6\linewidth}}{
        \cellcolor[rgb]{0.9,0.9,0.9}{
            \makecell[{{p{\linewidth}}}]{
                \texttt{\tiny{[P1$\langle$GM]}}
                \texttt{RIFIUTO: \{'A30', 'A47', 'A39', 'C24', 'C73', 'A19', 'B91', 'A75', 'A15', 'A56'\}} \\
\\ 
\texttt{ARGOMENTAZIONE: \{'Apprezzo che tu mantenga C24, C73 e A47. Tuttavia, la tua selezione usa molto sforzo su elementi medi (A15, A56, A75) e lascia fuori un moltiplicatore di valore per me. Posso rinunciare ad A39 se inseriamo B11: lo scambio A39→B11 aumenta il valore netto per me e restiamo entro il limite. Per restare flessibili, manteniamo A19 e B91 come pezzi leggeri/efficienti.'\}} \\
\\ 
\texttt{PROPOSTA: \{'A30', 'A47', 'C24', 'C73', 'A19', 'B91', 'B11', 'A15'\}} \\
            }
        }
    }
    & \\ \\

    \theutterance \stepcounter{utterance}  
    & \multicolumn{4}{p{0.6\linewidth}}{
        \cellcolor[rgb]{0.9,0.9,0.9}{
            \makecell[{{p{\linewidth}}}]{
                \texttt{\tiny{[P1$\rangle$GM]}}
                \texttt{RIFIUTO: \{'A30', 'A47', 'C24', 'C73', 'A19', 'B91', 'B11', 'A15'\}} \\
\\ 
\texttt{ARGOMENTAZIONE: \{'Posso accettare di lasciare fuori A39 per accomodare le tue preferenze, ma B11 a quel costo di sforzo ci forza a tagliare troppi pezzi che, combinati, rendono di più. Manteniamo i tuoi cardini C24, C73, A47. Ti propongo un compromesso ad alta efficienza: inseriamo A39 (che avevi offerto di sacrificare) solo se necessario, ma per ora puntiamo su un paniere ampio senza B11 che massimizza valore/sforzo e conserva i “leggeri” A19 e B91 come desideri.'\}} \\
\\ 
\texttt{PROPOSTA: \{'A30', 'A47', 'C24', 'C73', 'A19', 'B91', 'A75', 'A56', 'A15'\}} \\
            }
        }
    }
    & & \\ \\

    \theutterance \stepcounter{utterance}  
    & & \multicolumn{4}{p{0.6\linewidth}}{
        \cellcolor[rgb]{0.9,0.9,0.9}{
            \makecell[{{p{\linewidth}}}]{
                \texttt{\tiny{[GM$\rangle$P2]}}
                \texttt{RIFIUTO: \{'A30', 'A47', 'C24', 'C73', 'A19', 'B91', 'B11', 'A15'\}} \\
\\ 
\texttt{ARGOMENTAZIONE: \{'Posso accettare di lasciare fuori A39 per accomodare le tue preferenze, ma B11 a quel costo di sforzo ci forza a tagliare troppi pezzi che, combinati, rendono di più. Manteniamo i tuoi cardini C24, C73, A47. Ti propongo un compromesso ad alta efficienza: inseriamo A39 (che avevi offerto di sacrificare) solo se necessario, ma per ora puntiamo su un paniere ampio senza B11 che massimizza valore/sforzo e conserva i “leggeri” A19 e B91 come desideri.'\}} \\
\\ 
\texttt{PROPOSTA: \{'A30', 'A47', 'C24', 'C73', 'A19', 'B91', 'A75', 'A56', 'A15'\}} \\
            }
        }
    }
    & \\ \\

    \theutterance \stepcounter{utterance}  
    & & & \multicolumn{4}{p{0.6\linewidth}}{
        \cellcolor[rgb]{0.9,0.9,0.9}{
            \makecell[{{p{\linewidth}}}]{
                \texttt{\tiny{[GM$\langle$P2]}}
                \texttt{RIFIUTO: \{'A30', 'A47', 'C24', 'C73', 'A19', 'B91', 'A75', 'A56', 'A15'\}} \\
\\ 
\texttt{ARGOMENTAZIONE: \{'Capisco l’attenzione all’efficienza, e apprezzo che tu mantenga C24, C73 e A47 insieme ad A19 e B91. Tuttavia, con A15, A56 e A75 stiamo allocando molto sforzo su tre pezzi medi mentre lasciamo fuori un singolo pezzo per me estremamente prezioso. Propongo un compromesso: tolgo uno dei medi (A56) e reinserisco un pezzo forte per me (B11) mantenendo il paniere ampio e sotto limite. Questo conserva i tuoi cardini C24, C73, A47 e i “leggeri” A19, B91.'\}} \\
\\ 
\texttt{PROPOSTA: \{'A30', 'A47', 'C24', 'C73', 'A19', 'B91', 'A75', 'A15', 'B11'\}} \\
            }
        }
    }
     \\ \\

    \theutterance \stepcounter{utterance}  
    & & \multicolumn{4}{p{0.6\linewidth}}{
        \cellcolor[rgb]{0.9,0.9,0.9}{
            \makecell[{{p{\linewidth}}}]{
                \texttt{\tiny{[P1$\langle$GM]}}
                \texttt{RIFIUTO: \{'A30', 'A47', 'C24', 'C73', 'A19', 'B91', 'A75', 'A56', 'A15'\}} \\
\\ 
\texttt{ARGOMENTAZIONE: \{'Capisco l’attenzione all’efficienza, e apprezzo che tu mantenga C24, C73 e A47 insieme ad A19 e B91. Tuttavia, con A15, A56 e A75 stiamo allocando molto sforzo su tre pezzi medi mentre lasciamo fuori un singolo pezzo per me estremamente prezioso. Propongo un compromesso: tolgo uno dei medi (A56) e reinserisco un pezzo forte per me (B11) mantenendo il paniere ampio e sotto limite. Questo conserva i tuoi cardini C24, C73, A47 e i “leggeri” A19, B91.'\}} \\
\\ 
\texttt{PROPOSTA: \{'A30', 'A47', 'C24', 'C73', 'A19', 'B91', 'A75', 'A15', 'B11'\}} \\
            }
        }
    }
    & \\ \\

    \theutterance \stepcounter{utterance}  
    & \multicolumn{4}{p{0.6\linewidth}}{
        \cellcolor[rgb]{0.9,0.9,0.9}{
            \makecell[{{p{\linewidth}}}]{
                \texttt{\tiny{[P1$\rangle$GM]}}
                \texttt{RIFIUTO: \{'A30', 'A47', 'C24', 'C73', 'A19', 'B91', 'A75', 'A15', 'B11'\}} \\
\\ 
\texttt{ARGOMENTAZIONE: \{'Accolgo il taglio di A56, ma l’inserimento di B11 continua a comprimere troppo il resto del paniere per un ritorno complessivo inferiore. Manteniamo i tuoi cardini C24, C73 e A47, e i “leggeri” A19 e B91. Ti propongo un compromesso senza B11 che massimizza l’ampiezza e resta molto efficiente; se ti serve spazio per un tuo preferito, posso valutare lo scambio A75↔C75 oppure A15↔C26.'\}} \\
\\ 
\texttt{PROPOSTA: \{'A30', 'A47', 'C24', 'C73', 'A19', 'B91', 'A75', 'A15', 'A39'\}} \\
            }
        }
    }
    & & \\ \\

    \theutterance \stepcounter{utterance}  
    & & \multicolumn{4}{p{0.6\linewidth}}{
        \cellcolor[rgb]{0.9,0.9,0.9}{
            \makecell[{{p{\linewidth}}}]{
                \texttt{\tiny{[GM$\rangle$P2]}}
                \texttt{RIFIUTO: \{'A30', 'A47', 'C24', 'C73', 'A19', 'B91', 'A75', 'A15', 'B11'\}} \\
\\ 
\texttt{ARGOMENTAZIONE: \{'Accolgo il taglio di A56, ma l’inserimento di B11 continua a comprimere troppo il resto del paniere per un ritorno complessivo inferiore. Manteniamo i tuoi cardini C24, C73 e A47, e i “leggeri” A19 e B91. Ti propongo un compromesso senza B11 che massimizza l’ampiezza e resta molto efficiente; se ti serve spazio per un tuo preferito, posso valutare lo scambio A75↔C75 oppure A15↔C26.'\}} \\
\\ 
\texttt{PROPOSTA: \{'A30', 'A47', 'C24', 'C73', 'A19', 'B91', 'A75', 'A15', 'A39'\}} \\
            }
        }
    }
    & \\ \\

    \theutterance \stepcounter{utterance}  
    & & & \multicolumn{4}{p{0.6\linewidth}}{
        \cellcolor[rgb]{0.9,0.9,0.9}{
            \makecell[{{p{\linewidth}}}]{
                \texttt{\tiny{[GM$\langle$P2]}}
                \texttt{RIFIUTO: \{'A30', 'A47', 'C24', 'C73', 'A19', 'B91', 'A75', 'A15', 'A39'\}} \\
\\ 
\texttt{ARGOMENTAZIONE: \{'Apprezzo l’apertura agli scambi e il reinserimento di A39. Per trovare un vero compromesso senza B11, propongo di sostituire A75 con C75: lo sforzo resta simile ma per me il ritorno è migliore. Manteniamo i cardini condivisi (A47, C24, C73) e i leggeri (A19, B91). Se preferisci, posso anche cedere A15 in favore di C26 in una seconda iterazione, ma proverei prima con lo scambio A75→C75.'\}} \\
\\ 
\texttt{PROPOSTA: \{'A30', 'A47', 'C24', 'C73', 'A19', 'B91', 'C75', 'A15', 'A39'\}} \\
            }
        }
    }
     \\ \\

    \theutterance \stepcounter{utterance}  
    & & \multicolumn{4}{p{0.6\linewidth}}{
        \cellcolor[rgb]{0.9,0.9,0.9}{
            \makecell[{{p{\linewidth}}}]{
                \texttt{\tiny{[P1$\langle$GM]}}
                \texttt{RIFIUTO: \{'A30', 'A47', 'C24', 'C73', 'A19', 'B91', 'A75', 'A15', 'A39'\}} \\
\\ 
\texttt{ARGOMENTAZIONE: \{'Apprezzo l’apertura agli scambi e il reinserimento di A39. Per trovare un vero compromesso senza B11, propongo di sostituire A75 con C75: lo sforzo resta simile ma per me il ritorno è migliore. Manteniamo i cardini condivisi (A47, C24, C73) e i leggeri (A19, B91). Se preferisci, posso anche cedere A15 in favore di C26 in una seconda iterazione, ma proverei prima con lo scambio A75→C75.'\}} \\
\\ 
\texttt{PROPOSTA: \{'A30', 'A47', 'C24', 'C73', 'A19', 'B91', 'C75', 'A15', 'A39'\}} \\
            }
        }
    }
    & \\ \\

    \theutterance \stepcounter{utterance}  
    & \multicolumn{4}{p{0.6\linewidth}}{
        \cellcolor[rgb]{0.9,0.9,0.9}{
            \makecell[{{p{\linewidth}}}]{
                \texttt{\tiny{[P1$\rangle$GM]}}
                \texttt{RIFIUTO: \{'A30', 'A47', 'C24', 'C73', 'A19', 'B91', 'C75', 'A15', 'A39'\}} \\
\\ 
\texttt{ARGOMENTAZIONE: \{'Accolgo lo scambio A75→C75 come gesto verso le tue preferenze. Per completare il compromesso efficiente che hai suggerito, effettuo anche lo scambio A15→C26: il costo resta entro limite, manteniamo i tuoi cardini (A47, C24, C73) e i leggeri (A19, B91), e otteniamo un paniere ampio e bilanciato senza ricorrere a B11.'\}} \\
\\ 
\texttt{PROPOSTA: \{'A30', 'A47', 'C24', 'C73', 'A19', 'B91', 'C75', 'C26', 'A39'\}} \\
            }
        }
    }
    & & \\ \\

    \theutterance \stepcounter{utterance}  
    & & \multicolumn{4}{p{0.6\linewidth}}{
        \cellcolor[rgb]{0.9,0.9,0.9}{
            \makecell[{{p{\linewidth}}}]{
                \texttt{\tiny{[GM$\rangle$P2]}}
                \texttt{RIFIUTO: \{'A30', 'A47', 'C24', 'C73', 'A19', 'B91', 'C75', 'A15', 'A39'\}} \\
\\ 
\texttt{ARGOMENTAZIONE: \{'Accolgo lo scambio A75→C75 come gesto verso le tue preferenze. Per completare il compromesso efficiente che hai suggerito, effettuo anche lo scambio A15→C26: il costo resta entro limite, manteniamo i tuoi cardini (A47, C24, C73) e i leggeri (A19, B91), e otteniamo un paniere ampio e bilanciato senza ricorrere a B11.'\}} \\
\\ 
\texttt{PROPOSTA: \{'A30', 'A47', 'C24', 'C73', 'A19', 'B91', 'C75', 'C26', 'A39'\}} \\
            }
        }
    }
    & \\ \\

    \theutterance \stepcounter{utterance}  
    & & & \multicolumn{4}{p{0.6\linewidth}}{
        \cellcolor[rgb]{0.9,0.9,0.9}{
            \makecell[{{p{\linewidth}}}]{
                \texttt{\tiny{[GM$\langle$P2]}}
                \texttt{ACCORDO: \{'A30', 'A47', 'C24', 'C73', 'A19', 'B91', 'C75', 'C26', 'A39'\}} \\
\\ 
\texttt{ARGOMENTAZIONE: \{'Accetto questa selezione: rispetta il limite di sforzo, mantiene i cardini condivisi (A47, C24, C73) e conserva i leggeri A19 e B91. Inoltre include A30 e A39 che per me sono preziosi, insieme alle tue preferenze C75 e C26. È un compromesso ampio e bilanciato senza ricorrere a B11.'\}} \\
            }
        }
    }
     \\ \\

    \theutterance \stepcounter{utterance}  
    & & & \multicolumn{2}{p{0.3\linewidth}}{
        \cellcolor[rgb]{0.9,0.9,0.9}{
            \makecell[{{p{\linewidth}}}]{
                \texttt{\tiny{[GM$|$GM]}}
                \texttt{game successful} \\
            }
        }
    }
    & & \\ \\

    \theutterance \stepcounter{utterance}  
    & & & \multicolumn{2}{p{0.3\linewidth}}{
        \cellcolor[rgb]{0.9,0.9,0.9}{
            \makecell[{{p{\linewidth}}}]{
                \texttt{\tiny{[GM$|$GM]}}
                \texttt{end game} \\
            }
        }
    }
    & & \\ \\

\end{supertabular}
}

\end{document}
